%%%%%%%%%%%%%%%%%%%%%%%%%%%%%%%%%%%%%%%%%
% Journal Article
% LaTeX Template
% Version 2.0 (February 7, 2023)
%
% This template originates from:
% https://www.LaTeXTemplates.com
%
% Author:
% Vel (vel@latextemplates.com)
%
% License:
% CC BY-NC-SA 4.0 (https://creativecommons.org/licenses/by-nc-sa/4.0/)
%
% NOTE: The bibliography needs to be compiled using the biber engine.
%
%%%%%%%%%%%%%%%%%%%%%%%%%%%%%%%%%%%%%%%%%

%----------------------------------------------------------------------------------------
%	PACKAGES AND OTHER DOCUMENT CONFIGURATIONS
%----------------------------------------------------------------------------------------



\documentclass{article}
\usepackage[a4paper, total={6in, 8in}]{geometry}

\usepackage{booktabs} % Required for better horizontal rules in tables

\usepackage{array} % Required for manipulating table columns
\usepackage{caption} % Required for customizing captions
\usepackage{enumitem} % Required for list customization
\usepackage{hyperref} % Required for links
\usepackage{graphicx} % Required for including images


\usepackage{tikz}
\usetikzlibrary{shapes,snakes}
\usepackage{amsmath,amssymb}

\usepackage{xcolor}
\PassOptionsToPackage{svgnames}{xcolor}
\usepackage{tcolorbox}
\usepackage{lipsum}
\tcbuselibrary{skins,breakable}
\usetikzlibrary{shadings,shadows}



\definecolor{BgBiblio}{RGB}{255,103,103}
\definecolor{BgTitleBiblio}{RGB}{255,0,0}

\newenvironment{ImportantBiblio}[1]{%
    \tcolorbox[beamer,%
    noparskip,breakable,
    colback=BgBiblio,colframe=BgTitleBiblio,%
    colbacklower=red!75!green,%
    title=#1]}%
    {\endtcolorbox}


\definecolor{BgExp}{RGB}{139,161,255}
\definecolor{BgTitleExp}{RGB}{0,46,255}

\newenvironment{ImportantExp}[1]{%
    \tcolorbox[beamer,%
    noparskip,breakable,
    colback=BgExp,colframe=BgTitleExp,%
    colbacklower=red!75!green,%
    title=#1]}%
    {\endtcolorbox}


\definecolor{BgProtocole}{RGB}{170,170,170}
\definecolor{BgTitleProtocole}{RGB}{55,55,55}


\newenvironment{ProtocoleExp}[1]{%
    \tcolorbox[beamer,%
    noparskip,breakable,
    colback=BgProtocole,colframe=BgTitleProtocole,%
    colbacklower=red!75!green,%
    title=#1]}%
    {\endtcolorbox}


\definecolor{BgEtude}{RGB}{143,255,131}
\definecolor{BgTitleEtude}{RGB}{15,156,0}


\newenvironment{EtudeExp}[1]{%
    \tcolorbox[beamer,%
    noparskip,breakable,
    colback=BgEtude,colframe=BgTitleEtude,%
    colbacklower=red!75!green,%
    title=#1]}%
    {\endtcolorbox}


\title{Journal} % Article title, use manual lines breaks (\\) to beautify the layout


\author{Nathan Berrit}


\begin{document}

\maketitle % Output the title section

%----------------------------------------------------------------------------------------
%	ARTICLE CONTENTS
%----------------------------------------------------------------------------------------

\section*{Septembre 2024}

\subsection*{16/09}

\begin{EtudeExp}{Calibration FWHM : ND4-01 col2}
16/09 :
    \begin{itemize}
        \item Test en polarisation pour l'OPO
        \item Calibration de la gamme OPO.
    \end{itemize}

17/09 :
    \begin{itemize}
        \item Test en polar SHG 
        \item Calibration de la gamme SHG
    \end{itemize}
    
\end{EtudeExp}


\section*{Octobre 2024}


\begin{EtudeExp}{Calibration}

    Calibration plus fine permettant aussi d'avoir $\Delta z \sim 0$ à toute longueur d'onde en jouant sur la position de la lentille de sonde et de pompe.

    Pour ça, il faut un compromis, parce que changer la divergence des faisceau pour synchroniser les z fait qu'ils sont parfois moins bien focalisés. J'ai donc fait une étude du 21 au 30 Septembre pour quantifier quelle était la meilleur position de la lentille de pompe pour avoir à la fois les z synchro et garder un signal DTT proche de sa valeur max mesurée (donc un faisceau bien focalisé). j'ai conclut que $x_{Lp2} = 10.00$ mm était la meilleur pos pour ça, mais il se trouve qu'elle ne marche pas sur toute la gamme de l'OPO.

    Les études en z pos sont dans le dossier "Nanodisk/Lithographies/test zpos". Les mesures ont été faites sur un gros disque (check le cahier de manip pour sa taille), afin de ne pas avoir de problème de signal.
    
    Voilà un exemple :


    Optimisation avec $x_{Lp2} = 12.60$ mm, à 700 nm.

    \includegraphics[width=0.7\textwidth]{../../Experimental/Salle_particules_unique/Nanodisks/Lithographies/test_zpos_2609/optimisation_z_12-6mm.png}

    Même optimisation à 700nm, avec $x_{Lp2} = 10.00$ mm

    \includegraphics[width=0.7\textwidth]{../../Experimental/Salle_particules_unique/Nanodisks/Lithographies/test_zpos_2609/optimisation_z_10mm.png}


    On voit que le $\Delta_z = 0$ est atteint avec un meilleur DTT, donc une meilleure focalisation du faiseau, à la fois SHG et OPO.
\end{EtudeExp}


\begin{EtudeExp}{Conclusion calibration octobre 2024}
    Il y aura deux positions $x_{Lp2}$ : \begin{itemize}
        \item à 10.00mm pour utiliser l'OPO entre 590 et 650 nm : bonne configuration qui focalise au mieux la source SHG et la source OPO, tout en ayant $\Delta z = 0$.
        \item à 12.00mm pour utilise l'OPO entre 660 et 720 nm : moins bonne focalisation SHG, mais permet $\Delta z = 0$.
    \end{itemize}
    Les cartes à prendre sont : \begin{itemize}
        \item le 01 Oct : carte de 500 à 590 nm, avec $x_{Lp2} = 10.00$ mm.
        \item Le 04 Oct : carte de 610 à 650 nm, avec $x_{Lp2} = 10.00$ mm.
        \item Le 07 Oct : carte de 520 à 600 nm, avec $x_{Lp2} = 10.00$ mm.
        \item Le 07 Oct : carte de 660 à 720 nm, avec $x_{Lp2} = 12.00$ mm.
        \end{itemize}
\end{EtudeExp}

\section*{Decembre 2024}

\begin{EtudeExp}{Influence de la longueur d'onde de sonde pour l'étude acoustique}

    En Novembre et Décembre, j'ai fait l'étude acoustique de l'échantillon AR 2.3 de Sylvie.
    Les disques repérés par Sylvie étaient dans les zones rayées du substrat. Sylvie a fait ce choix parce que ça les rend repérables. Cependant, la présence de rayure rend l'environnement très bruit autour des ND. J'ai donc choisis de travailler à l'aveugle, en faisant des repérages sur des cartes SMS 20x20 um, et de sélectionner les objets candidats. 
    
    
\end{EtudeExp}





\section*{Janvier 2025}
\begin{ImportantBiblio}{Interpolation AR sur Python}


    Deux méthodes trouvées sur \href{https://stackoverflow.com/questions/37872171/how-can-i-perform-two-dimensional-interpolation-using-scipy}{ce lien} : \begin{itemize}
        \item griddata : on donne une série de points connus, et on demande une évaluation sur uné série de points
        \item Radial Basis Interpolation (RBF) : on donne une série de points connus et ça fait une fonction interpolée. 
    \end{itemize}

    En comparant mon programme d'interpolation 2D avec ce que renvoie wolfram mathematica. Quand on reste dans la zone balisée par l'interpolation, les résultats ont l'air proche et concordant (10 pourcent d'écart, pas plus). Hors de la zone, les résutats sont très différents et tout deux ont l'air de faire n'importe quoi. C'est flagrant pour la plupart de nos disques, qui ont une résonance à 570nm, qu'aucun de spectres simulée n'a pu montrer.



\end{ImportantBiblio}

\end{document}  